\documentclass[10pt,a4paper]{amsart}
\usepackage[margin=19mm]{geometry}
\usepackage{amsmath,amssymb,mathtools}
\usepackage[scr=rsfs,cal=euler]{mathalfa}
\usepackage{xfrac}
\usepackage[unicode,hidelinks,bookmarks=false]{hyperref}
\newcommand{\ie}{i.e.\ }
\newcommand{\cf}{cf.\ }
\newcommand{\resp}{respectively\ }
\newcommand{\Subsection}{\subsection}
\providecommand{\nobreakdash}{\nobreak\hbox{-}}
\setlength{\emergencystretch}{2em}
\multlinegap0pt
\usepackage{bm}
\usepackage{bbm}
\usepackage{dsfont}
\usepackage{xspace}
\usepackage{cancel}
\usepackage{mathtools}
\usepackage{amsthm}
\makeatletter
\@ifundefined{corollary}{\newtheorem{corollary}{Corollary}}{}
\@ifundefined{lemma}{\newtheorem{lemma}{Lemma}}{}
\makeatother
\newcommand\numberthis{\addtocounter{equation}{1}\tag{\theequation}}
\newcommand{\one}{\mathbf{1}}
\title[Optimal community detection in dense bipartite graphs]{Optimal community detection in dense bipartite graphs}
\author{Julien Chhor, Parker Knight}
\date{Source: arXiv:2505.18372v1 (CC BY 4.0)}
\begin{document}
\maketitle

\noindent\emph{Source and licence.} Optimal community detection in dense bipartite graphs. Version arXiv:2505.18372v1 (CC BY 4.0), distributed under the Creative Commons Attribution 4.0 International licence, as recorded in the arXiv licence metadata for this exact version and shown on the abstract page https://arxiv.org/abs/2505.18372v1. Full rights notice: licenses/arxiv-2505.18372-CC-BY-4.0.txt.\par\smallskip

\noindent\textit{Editorial scope.} Counted unit: the Lemma whose source label is \texttt{leq\_truncchisq\_concentration} in the article (source0.tex lines 2273--2281, source label \texttt{leq\_truncchisq\_concentration}), delivered together with the article's own complete original proof of it (source0.tex lines 2282--2445, one \texttt{proof} environment of 164 lines). No mathematics is written, repaired or re-derived: statement and proof are the article's own text line for line. Required context retained uncounted: lemma17\_LiuGaoSamworth (lines 2039--2047); eq\_def\_W\_appendix (lines 2258--2265); eq\_def\_gamma (lines 2258--2265); eq\_def\_nu\_a (lines 2258--2265); cor\_lemma17\_LiuGaoSamworth (lines 2268--2270); lem\_rescalen\_Bernstein (lines 2665--2673); lem\_Stirling (lines 2678--2684); lem\_W\_cvx\_increasing (lines 2698--2700); lem\_w\_quadratic (lines 2897--2904). Every definition, display and label of those lines is retained unchanged. Omitted from this package and not redistributed: 1--2038, 2048--2257, 2266--2267, 2271--2272, 2446--2664, 2674--2677, 2685--2697, 2701--2896, 2905--2959. Nothing inside a retained excerpt is omitted or altered: every retained body is delivered exactly as the article's lines are written, so no omission receipt of any kind accompanies this package. Citations: the retained excerpts carry no citation command at all, so no bibliography entry of the article is transcribed and no reference is redistributed. Reference adaptation: none was needed and none was made; every reference inside the delivered unit resolves inside it, so no unresolved reference or citation is produced. Printed slips kept literal: no printed word, symbol, hypothesis or number of the counted statement or of its proof was altered. Layout and class: the article's own class is \texttt{article} with packages including inputenc, fontenc, hyperref, url, booktabs, amsfonts, amsmath, amsthm, amssymb, nicefrac, microtype, xcolor, geometry, enumerate; the wrapper is the lane's pinned \texttt{amsart} editorial wrapper at 10pt on A4, which restates the theorem-like environments the retained text needs and adds the article's own macro definitions that the retained text needs (\texttt{\textbackslash numberthis}, \texttt{\textbackslash one}), with extra wrapper packages (bm, bbm, dsfont, xspace, cancel, mathtools). The delivered numbering is the unit's own. Assets: no retained span contains a figure environment, a graphics-inclusion command, a PDF-page inclusion, a table or a TikZ picture; no image file is copied into this package, the package declares no asset dependency and no image file is redistributed with it. Licence: version arXiv:2505.18372v1 of this article is distributed under the Creative Commons Attribution 4.0 International licence, as recorded in the arXiv licence metadata for this exact version and shown on the abstract page https://arxiv.org/abs/2505.18372v1. A full rights notice is delivered at licenses/arxiv-2505.18372-CC-BY-4.0.txt. Characters: every retained excerpt is pure ASCII, so the delivered engine, XeLaTeX with its default Latin Modern text font, reports no missing character.
\begin{lemma}\label{lemma17_LiuGaoSamworth}
Let $n \in \mathbb N$ and $p \in (0,1/4)$. 
Let $X \sim \operatorname{Bin}(n,p)$ and, for any $a>0$,  define 
\begin{align*}
    &Z \,=\, (X-np)/\sigma, \qquad \text{ where } \sigma = \sqrt{np(1-p)}\\
    &v(x)  = \sigma x \log\left(1+ \frac{x}{\sigma}\right), \quad \forall x > 0.
\end{align*}
For any $\alpha\geq 1$, there exist two constants $C_\alpha, \bar C>0$ such that, for any $a \in\big[C_\alpha, \sigma\big]$, we have $\mathbb E\left[v(Z)^{\alpha} \big|Z\geq a\right] \leq \bar C v(a)^{\alpha}$. %It follows that $\nu_a \leq \bar C a^2$ and $\gamma_a \leq \bar C a^4$ for some absolute constant $\bar C$. 
\end{lemma}

\begin{align*}
    w(x) &= n(1-p_0) \, h_B\!\left(-\frac{x-np_0}{n(1-p_0)}\right) + np_0 \, h_B\!\left(\frac{x-np_0}{np_0}\right), \forall x \in [0, n]\numberthis \label{eq_def_W_appendix}\\
    Z &= (Y - np_0)/\sigma, \qquad \text{ where } \sigma = \sqrt{np_0(1-p_0)}\\[5pt]
    W &= w(Y)\\
    \nu_a &=  \mathbb 
    E\left[W| Z \geq a\right]\numberthis \label{eq_def_nu_a}\\
\gamma_a &= \mathbb E\left[W^2 | Z \geq a\right].\numberthis\label{eq_def_gamma}
\end{align*}

\begin{corollary}\label{cor_lemma17_LiuGaoSamworth}
    Recall the definitions of the functions $\nu$ and $\gamma$ from~\eqref{eq_def_nu_a} and~\eqref{eq_def_gamma}. There exist sufficiently large constants $C, \bar C$ such that, for any $a \in [C, \sigma]$, it holds that $\nu_a \leq \bar C \sigma a \log\left(1+\frac{a}{\sigma}\right) \leq \bar C a^2$ and $\gamma_a \leq \bar C \sigma a^2 \log^2\left(1+\frac{a}{\sigma}\right) \leq \bar C a^4$.
\end{corollary}

\begin{lemma}\label{lem_rescalen_Bernstein}
    Let $n\in \mathbb N$ and $p \in (0,1/2)$. 
    Let $X \sim \operatorname{Bin}(n,p)$ and $\sigma = \sqrt{np(1-p)}$.
    For any $t > 0$, we have
    \begin{align*}
        \mathbb P\left(\frac{X - np}{\sigma} \geq t\right) \leq \exp\left(- \frac{t^2/2}{1+\frac{t}{3\sigma}}\right)\\
        \mathbb P\left(\frac{X - np}{\sigma} \leq -t\right) \leq \exp\left(- \frac{t^2/2}{1+\frac{t}{3\sigma}}\right)
    \end{align*}
\end{lemma}

\begin{lemma}\label{lem_Stirling}
    There exist two absolute constants $C_0, c_0$ such that, for any $n,k \in \mathbb N$ such that $k < n$ and $p \in (0,1)$, we have
    \begin{align*}
        {n \choose k} p^k (1-p)^{n-k} &\leq \frac{C_0}{c_0^2\sqrt{2\pi}} \sqrt{\frac{n}{(n-k)k}} \left(\frac{n(1-p)}{n-k}\right)^{n-k} \left(\frac{np}{k}\right)^k\\
        & = \frac{C_0}{c_0^2\sqrt{2\pi}} \sqrt{\frac{n}{(n-k)k}} \exp\left(-n(1-p) h_B\left(-\frac{k-np}{n(1-p)}\right) - np \, h_B\left(\frac{k-np}{np}\right)\right). 
    \end{align*}
\end{lemma}

\begin{lemma}\label{lem_W_cvx_increasing}
    The function $w$ defined in~\eqref{eq_def_W_appendix} is convex and increasing over $[np_0, n)$.
\end{lemma}

\begin{lemma}\label{lem_w_quadratic}
    Let $n\in \mathbb N$ and $p_0 \in (0,1)$ and $w$ denote the function defined in~\eqref{eq_def_W_appendix}. For any $x \in [0,n]$, the following properties hold.
    \begin{enumerate}
        \item We have $w(x) \leq \frac{(x-np_0)^2}{ \sigma^2}$ where $\sigma^2 = np_0(1-p_0)$
        \item There exists a universal constant $c_0>0$ such that, for any $x \in [np_0, np_0 + c_0\sigma^2]$, we have $w(x) \geq \frac{(x-np_0)^2}{8 \sigma^2}$.
        \item For any $x \in [c_0 \sigma^2, n]$, we have $w(x) \asymp \sigma \frac{x-np_0}{\sigma} \log\left(1+ \frac{x-np_0}{\sigma^2}\right).$
    \end{enumerate}
\end{lemma}

\begin{lemma}\label{leq_truncchisq_concentration} Let $m,n \in \mathbb N$ and $p_0 \in (0,1/4)$. 
Let $Y_1,\dots, Y_m \overset{iid}{\sim} \operatorname{Bin}(n,p_0)$ and for any $j \in \{1,\dots,m\}$ define 
$Z_j = (Y_j-np_0)/\sigma, \text{ where } \sigma = \sqrt{np_0(1-p_0)}.$ 
For any $j \in \{1,\dots,m\}$, define $W_j = w(Y_j)$ where $w$ is defined in~\eqref{eq_def_W_appendix} and recall the definition of $\nu_a$ from~\eqref{eq_def_nu_a} for $a > 0$. 
Then there exist universal constants $C,C^{*}, c, c'>0$ such that for any $a\in [C,c\sigma]$ and $x>0$, we have
$$
\mathbb{P}\left(\sum_{j=1}^{m}\left(W_j-\nu_{a}\right) \one_{\left\{Z_{j} \geq a\right\}} \geq C^{*}\left(\sqrt{m e^{-c'a^2} x}+x\right)\right) \leq e^{-x}.
$$
\end{lemma}

\begin{proof}[Proof of Lemma~\ref{leq_truncchisq_concentration}]
In this proof, we will write $p$ rather than $p_0$ to alleviate the notation.  
Let $Y \sim \operatorname{Bin}(n,p)$ and $Z = (Y-np)/\sigma$, and consider $X=\left(W-\nu_{a}\right) \one_{\{Z \geq a\}}$ where
\begin{align*}
    W = n(1-p)\,h_B\!\left(-\frac{Y-np}{n(1-p)}\right) + np \,h_B\!\left(\frac{Y-np}{np}\right).
\end{align*}
We first derive a bound for the moment-generating function $\mathbb{E}\left(e^{\lambda X}\right)$. Since $\mathbb{E}(X)=0$, we have
$$
\mathbb{E}\left(e^{\lambda X}\right)=1+\mathbb{E}\left(e^{\lambda X}-1-\lambda X\right).
$$

By the deterministic bound
$$
e^{x}-1-x \leq \begin{cases}(-x) \wedge x^{2}, & \text { if } x<0 \\ x^{2}, & \text { if } 0 \leq x \leq 1 \\ e^{x}, & \text { if } x \geq 1\end{cases}
$$
we have
$$
\mathbb{E}\left(e^{\lambda X}-1-\lambda X\right) \leq \lambda^{2} \mathbb{E}\left(X^{2} \one_{\{X<0\}}\right)+\lambda^{2} \mathbb{E}\left(X^{2} \one{\{0 \leq X \leq 1 / \lambda\}}\right)+\mathbb{E}\left(e^{\lambda X} \one_{\{X\geq1 / \lambda\}}\right)
$$
and we will bound the three terms separately. By Corollary~\ref{cor_lemma17_LiuGaoSamworth}, the first term is bounded as
\begin{align*}
    \mathbb{E}\left(\big(|X|\land X^{2}\big) \one_{\{X<0\}}\right) & \leq \mathbb{P}(X<0)\Big[\big|\nu_{a}-w(np+a\sigma)\big|\land \big(\nu_{a}-w(np+a\sigma)\big)^{2}  \Big]\\
    & \leq \mathbb{P}(X<0)\Big[\big|\nu_{a}-\frac{a^2}{8}\big|\land \big(\nu_{a}-\frac{a^2}{8}\big)^{2}  \Big] \qquad \text{ by Lemma~\ref{lem_w_quadratic}.2}\\[5pt]
    & \leq \mathbb{P}(Z \geq a)\left(\nu_{a}-\frac{a^{2}}{8}\right)^{2} \\[5pt]
    &\leq (\bar C-1/8)^2 a^4 \mathbb P\left(Y-np \geq a \sigma\right) \qquad \text{ by Corollary~\ref{cor_lemma17_LiuGaoSamworth}}\\
    & \leq (\bar C-1/8)^2 a^4\exp \left(-\frac{a^2/2}{1 + \frac{a/\sigma}{3}}\right)\qquad \text{ by Lemma~\ref{lem_rescalen_Bernstein}}\\
    & = (\bar C-1/8)^2 a^4\exp \left(-\frac{3a^2}{8}\right)\\
    & \leq \exp\left(- \frac{a^2}{4}\right)
\end{align*}
provided the universal constant $C$ such that $a \geq C$ is chosen large enough depending on $\bar C$. We now bound the second term. Recalling the definition of the function $\gamma$ from~\eqref{eq_def_gamma}, we have
$$
\begin{aligned}
    \mathbb{E}\left(X^{2} \one_{\{0<X \leq 1 / \lambda\}}\right)&\leq \mathbb{E}\left(X^{2} \one_{\{X>0\}}\right) = \mathbb E\left[\left(W-\nu_a\right)^2 \one_{\{Z >\sqrt{\nu_a}\}}\right]\\[5pt]
    & \leq 2 \mathbb E\left[W^2 \one_{\{Z >\sqrt{\nu_a}\}}\right] + 2 \nu_a^2 \mathbb P\left(Z>\sqrt{\nu_a}\right)\\[10pt]
    & = 2 \left(\gamma_{\sqrt{\nu_a}} + \nu_a^2\right) \mathbb P\left(Z>\sqrt{\nu_a}\right) \\
    & \leq 2\left(\bar C \nu_a^2 + \nu_a^2\right) \exp\left(- \frac{\nu_a/2}{1+\frac{\sqrt{\nu_a}}{3\sigma}}\right) \qquad \text{ by Lemmas~\ref{lemma17_LiuGaoSamworth} and~\ref{lem_rescalen_Bernstein} }\\
    & \leq 2 \left(\bar C^3 + \bar C^2\right)a^4 \exp\left(-\frac{\nu_a/2}{1+ \sqrt{\bar C}/3}\right) \quad \text{ by Lemma~\ref{lemma17_LiuGaoSamworth} and } a \leq \sigma\\
    & \leq \exp\left(-\frac{\nu_a}{2\bar C}\right)\\
    & \leq \exp\left(-\frac{a^2}{2\bar C}\right)
\end{aligned}
$$
provided the universal constant $C$ such that $a \geq C$ is chosen large enough depending on $\bar C$.

Finally we control the third term. 
We note that the event $X\geq\frac{1}{\lambda}$ is equivalent to $W \geq \nu_a + \frac{1}{\lambda}$. 
We let $k_0 \in \mathbb N$ be the unique integer larger than or equal to $np$ such that the event $\big\{W \geq \nu_a + \frac{1}{\lambda}$ and $Z \geq 0 \big\}$ is equivalent to $\{Y \geq k_0\}$. 
This integer exists by Lemma~\ref{lem_W_cvx_increasing} since $w$ is increasing over $[np,n]$ by Lemma~\ref{lem_W_cvx_increasing}, $w(np + a\sigma) = 0$, $\lim_{x \to n} w(x) = 2 n \log(1/p) > \nu_a + \frac{1}{\lambda}$ for $\lambda \leq \frac{1}{2}$ and $\nu_a \leq \bar C a^2 \leq \bar C c^2 \sigma^2 \leq \bar C c n < 2 n \log(1/p)$ for $c$ chosen sufficiently small depending on $\bar C$. 
Moreover, we also have $k_0 \geq np + a\sigma$ since 
$w(np+a\sigma) \leq a^2 < \nu_a + \frac{1}{\lambda}$.
Now, we have
\begin{align*}
\mathbb E\left[e^{\lambda X} \one{\big(X\geq1/\lambda\big)}\right] &= \mathbb E\left[\exp\left(\lambda(W - \nu_a)\right)\one{\big(W \geq \nu_a + \frac{1}{\lambda}\big)}\right]\\
& = e^{-\lambda \nu_a}  \sum_{k = k_0}^n e^{\lambda w(k)}\mathbb P(Y = k)\\
& \leq e^{-\lambda \nu_a}   e^{\lambda w(n)}\mathbb P(Y = n) + e^{-\lambda \nu_a}   e^{\lambda w(k_0)}\mathbb P(Y = k_0) \\
& \qquad + \frac{C_0}{\sqrt{2\pi}c_0^2} e^{-\lambda \nu_a} \sum_{k = k_0}^{n-1} \sqrt{\frac{n}{(n-k)k}} \\
& \qquad \times \exp\left((2\lambda - 1) \left\{n(1-p) h_B\left(-\frac{k-np}{n(1-p)}\right) + np \, h_B\left(\frac{k-np}{np}\right)\right\}\right) \\
&\hspace{3cm} \text{ (by Lemma~\ref{lem_Stirling})}\\
& = e^{-\lambda \nu_a}   e^{\lambda w(n)}\mathbb P(Y = n) +  e^{-\lambda \nu_a}   e^{\lambda w(k_0)}\mathbb P(Y = k_0) \\
& \qquad + \frac{C_0}{\sqrt{2\pi}c_0^2} e^{-\lambda \nu_a} \sum_{I_1 \cup I_2} \sqrt{\frac{n}{(n-k)k}} e^{\left(\lambda - \frac{1}{2}\right)w(k)}\numberthis \label{eq_thirn_term_moment_generating_function}
\end{align*}
where we have defined $I_1 = \{k_0+1, \dots, \lfloor C^* \sigma^2\rfloor \}$ and $I_2 = \{\lfloor C^* \sigma^2\rfloor + 1, \dots, n-1\}$ for some constant $C^*$ that will be adjusted later. 
We first compute the sum over the set of indices $I_1$. 
Below, we use the notation $C'$ to denote a constant whose value may change in each appearance. 
Using the change of variables 
\begin{align*}
    & y = w(x)\\
    & dx = \frac{dy}{w' \circ w^{-1}(y)} = \left(\log\left(\frac{w^{-1}(y)}{np} \frac{n(1-p)}{n-w^{-1}(y)}\right)\right)^{-1} dy
\end{align*}
we have  
\begin{align*}
    & \quad \sum_{I_1} \sqrt{\frac{n}{(n-k)k}} e^{\left(\lambda - \frac{1}{2}\right)w(k)}\\
    & \leq C'\int_{k_0}^{\lfloor C^*\sigma\rfloor +1} \sqrt{\frac{n}{(n-x) x}} e^{(\lambda - \frac{1}{2}) w(x)} dx\\
    & \leq \frac{C'}{\sigma}\int_{k_0}^{\lfloor C^*\sigma\rfloor +1}  e^{(\lambda - \frac{1}{2}) w(x)} dx\\
    & = \frac{C'}{\sigma} \int_{w(k_0)}^{w(\lfloor C^*\sigma^2\rfloor + 1)} e^{(\lambda - \frac{1}{2}) y} \left(\log\left(\frac{w^{-1}(y)}{np} \frac{n(1-p)}{n-w^{-1}(y)}\right)\right)^{-1} dy\\  
    & = \frac{C'}{\sigma} \int_{w(k_0)}^{w(\lfloor C^*\sigma^2\rfloor + 1)} e^{(\lambda - \frac{1}{2}) y} \left(\log\left(\frac{w^{-1}(y)}{np} \frac{1}{1 - \frac{w^{-1}(y) - np}{n(1-p)}}\right)\right)^{-1} dy\\
    & \leq \frac{C'}{\sigma} \int_{w(k_0)}^{w(\lfloor C^*\sigma^2\rfloor + 1)} e^{(\lambda - \frac{1}{2}) y} \left(\log\left(\frac{w^{-1}(y)}{np} \left(1+\frac{ w^{-1}(y) - np}{n(1-p)}\right)\right)\right)^{-1} dy.
\end{align*}
Moreover, we have 
$w(x) \leq \frac{(x-np)^2}{2\sigma^2}$ for any $x \in [np,n]$. 
Therefore, $w^{-1}(y) \geq 
 np + \sigma \sqrt{2y}$. Therefore, we can further control the display above as
 \begin{align*}
     & \quad \sum_{I_1} \sqrt{\frac{n}{(n-k)k}} e^{\left(\lambda - \frac{1}{2}\right)w(k)}\\
     & \leq \frac{C'}{\sigma} \int_{w(k_0)}^{w(\lfloor C^*\sigma^2\rfloor + 1)} e^{(\lambda - \frac{1}{2}) y} \left(\log\left(\left(1+ \frac{\sigma\sqrt{2y}}{np}\right) \left(1+\frac{\sigma\sqrt{2y}}{n(1-p)}\right)\right)\right)^{-1} dy\\
     & \leq \frac{C'}{\sigma} \int_{w(k_0)}^{w(\lfloor C^*\sigma^2\rfloor + 1)} e^{(\lambda - \frac{1}{2}) y} \left(\log\left(1+ \frac{\sqrt{2y}}{\sigma}\right)\right)^{-1} dy.
 \end{align*}
Moreover, we have $w(x) \leq \frac{(x-np)^2}{\sigma^2} $ for any $x >0$ by Lemma~\ref{lem_w_quadratic} so that $w(\lfloor C^*\sigma^2\rfloor + 1) \leq C' \sigma^2$, which yields that, for some constant $c'$ that depends on $C^*$, we have
\begin{align*}
    \log\left(1+ \frac{\sqrt{2y}}{\sigma}\right) \geq c'\frac{\sqrt{y}}{\sigma}, \qquad \forall y \in \left[{w(k_0)},{w(\lfloor C^*\sigma^2\rfloor + 1)} \right].
\end{align*}
We obtain
\begin{align*}
    \sum_{I_1} \sqrt{\frac{n}{(n-k)k}} e^{\left(\lambda - \frac{1}{2}\right)w(k)}
    & \leq \frac{C'}{\sigma} \int_{w(k_0)}^{w(\lfloor C^*\sigma^2\rfloor + 1)} e^{(\lambda - \frac{1}{2}) y} \frac{\sigma}{\sqrt{y}} dy\\[5pt]
    & \leq C' \int_{(\frac{1}{2} - \lambda)w(k_0)}^\infty e^{-z} \frac{1}{\sqrt{z}} \frac{1}{\sqrt{1/2 - \lambda}} dz\\[5pt]
    & \leq C' \int_{(\frac{1}{2} - \lambda)(\nu_a + \frac{1}{\lambda})}^\infty e^{-z} \frac{1}{\sqrt{z}} \frac{1}{\sqrt{1/2 - \lambda}} dz\\[5pt]
    & \leq \frac{C' \exp\left(- (\frac{1}{2} - \lambda)(\nu_a + \frac{1}{\lambda})\right)}{\sqrt{(\frac{1}{2} - \lambda)(\nu_a + \frac{1}{\lambda})}},
\end{align*}
which yields
\begin{align}
    \frac{C_0}{\sqrt{2\pi}c_0^2} e^{-\lambda \nu_a} \sum_{I_1} \sqrt{\frac{n}{(n-k)k}} e^{\left(\lambda - \frac{1}{2}\right)w(k)} \leq \frac{C_0}{\sqrt{2\pi}c_0^2} e^{-\lambda \nu_a} \frac{C' \exp\left(- (\frac{1}{2} - \lambda)(\nu_a + \frac{1}{\lambda})\right)}{\sqrt{(\frac{1}{2} - \lambda)(\nu_a + \frac{1}{\lambda})}}\label{eq_thirn_term_I1}
\end{align}


We now show that, for any $k \in I_2 \setminus \{n-1\}$, we have
\begin{align*}
    \sqrt{\frac{n}{(n-k-1)(k+1)}} e^{\left(\lambda - \frac{1}{2}\right)w(k+1)} \leq \sqrt{2} e^{-2} \sqrt{\frac{n}{(n-k)k}} e^{\left(\lambda - \frac{1}{2}\right)w(k)}.
\end{align*}
Indeed, we have
\begin{align*}
    \sqrt{\frac{n}{(n-k-1)(k+1)}}  \leq \sqrt{2}\sqrt{\frac{n}{(n-k)k}}.  
\end{align*}
Moreover, since the function $w$ is increasing and convex over $[np,n)$ by Lemma~\ref{lem_W_cvx_increasing}, we have 
\begin{align*}
    w(k+1) - w(k)\geq w'(k) &\geq  2 \log\left(\frac{k}{np} \frac{n(1-p)}{n-k}\right)\\
    & \geq 2 \log\left(\frac{k_0}{np} \frac{n(1-p)}{n-k_0}\right) \qquad \text{ since $w'$ is increasing over }[np,n]\\
    & \geq \log\left(\frac{np + \lfloor C^*\sigma^2 \rfloor}{np} \frac{n(1-p)}{n-np - \lfloor C^*\sigma^2 \rfloor}\right)\\
    & = \log\left(\left(1 + \frac{\lfloor C^*\sigma^2 \rfloor}{np} \right)\frac{1}{1 - \frac{\lfloor C^*\sigma^2 \rfloor}{n (1-p)}}\right)\\
    & \geq \log\left(\left(1 + \frac{\lfloor C^*\sigma^2 \rfloor}{np} \right)\left(1 + \frac{\lfloor C^*\sigma^2 \rfloor}{n (1-p)}\right)\right)\\
    & \geq \log \left(1+\frac{\lfloor C^*\sigma^2 \rfloor}{\sigma^2}\right)\\
    & \geq 8
\end{align*}
for $C^*$ larger than an absolute constant. 
For $\lambda \leq 1/4$, we therefore obtain
\begin{align*}
    \sqrt{\frac{n}{(n-k-1)(k+1)}} e^{\left(\lambda - \frac{1}{2}\right)w(k+1)} \leq \sqrt{2} e^{-2} \sqrt{\frac{n}{(n-k)k}} e^{\left(\lambda - \frac{1}{2}\right)w(k)}.
\end{align*}
Therefore, the sum over $I_2$ can be controlled as by a geometric series, and we obtain 
\begin{align*}
    \frac{C_0}{\sqrt{2\pi}c_0^2} e^{-\lambda \nu_a} \sum_{k \in I_2} \sqrt{\frac{n}{(n-k)k}} e^{\left(\lambda - \frac{1}{2}\right)w(k)} & \leq C' e^{-\lambda \nu_a} \sqrt{\frac{n}{(n-\lfloor C^* \sigma^2 \rfloor)\lfloor C^* \sigma^2 \rfloor}} e^{\left(\lambda - \frac{1}{2}\right)w(\lfloor C^* \sigma^2 \rfloor)}\\
    & \leq C'' e^{-\lambda \nu_a} \frac{1}{\sigma} e^{\left(\lambda - \frac{1}{2}\right)\left(\nu_a + 1/\lambda\right)}\numberthis\label{eq_thirn_term_I2}
\end{align*}
provided $C^*$ is large enough. 

To conclude, it remains to control the terms $e^{-\lambda \nu_a}   e^{\lambda w(n)}\mathbb P(Y = n) $ and $ e^{-\lambda \nu_a}   e^{\lambda w(k_0)}\mathbb P(Y = k_0)$. 
For $k \in \{k_0,n\}$, we have
\begin{align*}
    & \quad e^{-\lambda \nu_a}   e^{\lambda w(k)}\mathbb P(Y = k) \\
    & \leq e^{-\lambda \nu_a}   e^{\lambda w(k)}\frac{C_0}{c_0^2\sqrt{2\pi}} \sqrt{\frac{n}{(n-k)k}} \exp\left(-n(1-p) \,h_B\!\left(-\frac{k-np}{n(1-p)}\right) - np \, \, h_B\!\left(\frac{k-np}{np}\right)\right) \text{ by Lemma~\ref{lem_Stirling}}\\
    & \leq \frac{C_0}{c_0^2\sqrt{2\pi}} e^{-\lambda \nu_a}\sqrt{\frac{2}{k}} e^{(\lambda - \frac{1}{2}) w(k)}\\
    & \leq \frac{C_0}{c_0^2\sqrt{2\pi}} e^{-\lambda \nu_a}\sqrt{\frac{2}{\sigma^2}} e^{(\lambda - \frac{1}{2}) (\nu_a + \frac{1}{\lambda})}\numberthis \label{eq_thirn_term_remainder}
\end{align*}
by definition of $k_0$. 
Combining equations~\eqref{eq_thirn_term_moment_generating_function}, \eqref{eq_thirn_term_I1}, \eqref{eq_thirn_term_I2} and~\eqref{eq_thirn_term_remainder}, we obtain
\begin{align*}
    \mathbb E\left[e^{\lambda X} \one{\big(X\geq1/\lambda\big)}\right] &\leq C e^{-\lambda \nu_a + (\lambda - \frac{1}{2})(\nu_a + \frac{1}{\lambda})}\\
    & \leq C e^{-\frac{\nu_a}{2} - \frac{1}{2\lambda}} \\
    & \leq C \lambda^2 e^{- \frac{\nu_a}{2}}.
\end{align*}



Assembling these calculations, we have for $\lambda \in [0,M]$
\[\mathbb{E} e^{\lambda X} \leq 1+C' \lambda^{2} e^{- ca^{2}} \leq \exp \left(C' \lambda^{2} e^{-ca^{2}}\right)\]
where $C' > 0$ is a constant that depends on $\bar{C}$. Then, for any $t>0$, we have
$$
\begin{aligned}
\mathbb{P}\left(\sum_{j=1}^{m}\left(W_{j}-\nu_{a}\right) \one_{\left\{Z_{j} \geq a\right\}}>t\right) & \leq \inf _{\lambda \leq M} e^{-\lambda t}\left(\mathbb{E} e^{\lambda X}\right)^{m} \leq \inf _{\lambda \leq M} \exp \left(-\lambda t + C' m \lambda^{2} e^{-ca^{2}}\right) \\
& =\exp \left\{-\left(\frac{Mt^{2} e^{ca^{2}}}{C' m} \wedge \frac{t}{2M}\right)\right\}
\end{aligned}
$$

Setting $t=C^*\left(\sqrt{m e^{-ca^{2}} x}+x\right)$ for $C^*$ sufficiently large, we obtain the desired result.
\end{proof}

\end{document}
